\documentclass{article}

\usepackage{amsmath,amssymb}
\usepackage{xurl}
\usepackage[hidelinks]{hyperref}
\setlength{\emergencystretch}{2em}

% Shared commands for notes in docs/paper-gaps/.

\DeclareUnicodeCharacter{2080}{\ifmmode{}_{0}\else\textsubscript{0}\fi}
\DeclareUnicodeCharacter{2081}{\ifmmode{}_{1}\else\textsubscript{1}\fi}
\DeclareUnicodeCharacter{2082}{\ifmmode{}_{2}\else\textsubscript{2}\fi}
\DeclareUnicodeCharacter{2099}{\ifmmode{}_{n}\else\textsubscript{n}\fi}

\newcommand{\Lean}{\textsc{Lean}}
\ExplSyntaxOn
\NewDocumentCommand{\leanid}{m}{
  \ifmmode
    \textsf{\detokenize{#1}}
  \else
    \group_begin:
    \urlstyle{sf}
    \tl_set:Nn \l_tmpa_tl {#1}
    \tl_replace_all:Nnn \l_tmpa_tl {\_} {_}
    \exp_args:NV \path \l_tmpa_tl
    \group_end:
  \fi
}
\ExplSyntaxOff
\providecommand{\tr}{\operatorname{tr}}
\newcommand{\tnleanIssueBase}{https://github.com/LionSR/TNLean/issues/}
\newcommand{\tnleanPrBase}{https://github.com/LionSR/TNLean/pull/}
\newcommand{\tnleanDocBase}{https://sirui-lu.com/TNLean/docs/}
\newcommand{\ghissue}[1]{\href{\tnleanIssueBase#1}{GitHub issue \##1}}
\newcommand{\ghpr}[1]{\href{\tnleanPrBase#1}{PR \##1}}
\newcommand{\doclink}[1]{\href{\tnleanDocBase#1.html}{\path{#1}}}

% Dirac notation and the identity operator, as in blueprint/src/macros/common.tex.
% \providecommand keeps note-local definitions authoritative.
\usepackage{dsfont}
\providecommand{\ket}[1]{|#1\rangle}
\providecommand{\bra}[1]{\langle #1|}
\providecommand{\braket}[2]{\langle #1 | #2 \rangle}
\providecommand{\ketbra}[2]{|#1\rangle\!\langle #2|}
\providecommand{\Id}{\mathds{1}}
\providecommand{\one}{\mathds{1}}
\providecommand{\C}{\mathbb{C}}
\providecommand{\MN}[1]{M_{#1}(\C)}
\providecommand{\MD}{M_D(\C)}

% Verdict marker: \gapnote{<kind>}{<status>}. Kind is the policy
% classification (clarification, local-correction, scope-restriction,
% unfaithful, false-source, open-gap); status is open, wip (elimination
% underway), resolved, or historical. Machine-read by the site index;
% prints a verdict line.
\newcommand{\gapnote}[2]{\par\noindent\textbf{Verdict:} #1 (#2).\par}


% Fallbacks keep this note compilable when a different command.tex is found first.
\providecommand{\leanid}[1]{\textsf{\detokenize{#1}}}
\providecommand{\gapnote}[2]{\par\noindent\textbf{Verdict:} #1 (#2).\par}

\title{The Ground Energy in the Number of Sampled Words of Landau, Vazirani and Vidick}
\author{}
\date{2026-10-29}

\begin{document}
\maketitle
\gapnote{false-source}{open}

\section{The source statement}

Landau, Vazirani and Vidick~\cite{Landau2015Polynomial},
\href{https://arxiv.org/abs/1307.5143v1}{arXiv:1307.5143v1}, source
\path{algo53.tex}, consider $H=\sum_{i=1}^{n-1}H_i$ on $n$ qudits of dimension
$d$, with $H_i$ acting on the qudits $i$ and $i+1$ and $0\le H_i\le1$, a unique
ground state $\Gamma$ of energy $\varepsilon_0$, and gap
$\varepsilon=\varepsilon_1-\varepsilon_0$ (lines 283--285). No bound on
$\varepsilon_0$ is assumed. With $C=(1-\varepsilon_0/n)^{-1}$ and $P_i=1-H_i$,
the operator
\begin{align*}
    A=\Bigl(C\Bigl(1-\frac Hn\Bigr)\Bigr)^m
    =C^m\frac1{n^m}\sum_{I\in\{1,\ldots,n\}^m}P_I
\end{align*}
is approximated by the sampling operator
$K=C^m\ell^{-1}\sum_{j=1}^\ell P_{I_j}$ of $\ell$ independent uniform words
(Definition \texttt{d:sagsp}, line 611). Corollary \texttt{c:agsp} (lines
626--638) asserts that some $m=O(\varepsilon^{-1}n\log(q(n)/c_\varepsilon))$
and $\ell=n^{O(1/\varepsilon)}$ give $\|K-A\|\le1/q(n)$, with every letter
occurring at most $\kappa\log n$ times in every word, $\kappa=O(1/\varepsilon)$,
with probability at least $1-1/n^3$. The proof (lines 640--644) applies the
matrix Bernstein inequality with $R=C^m+1$ and $\sigma^2\le(C^m+1)^2$, and then
uses
\begin{align}
    C^{2m}=n^{O(1/\varepsilon)}.
    \label{eq:printed}
\end{align}
Theorem~1 of the source (\texttt{thm:main}, lines 213--216) states the running time
$n^{c(d,\varepsilon)}\,\mathrm{poly}(n/\eta)$ with
$c(d,\varepsilon)=2^{O(\log^3d/\varepsilon)}$, independent of $\varepsilon_0$.

\section{The point at issue}

Since $1+x\le e^x$,
\begin{align*}
    C=1+\frac{\varepsilon_0}{n-\varepsilon_0}
    \le\exp\Bigl(\frac{\varepsilon_0}{n-\varepsilon_0}\Bigr),
    \qquad
    C^m\le\exp\Bigl(\frac{\varepsilon_0m}{n-\varepsilon_0}\Bigr).
\end{align*}
At $m$ of order $(n/\varepsilon)\log n$ this gives
$C^{2m}=n^{O(\varepsilon_0/\varepsilon)}$, which is (\ref{eq:printed}) only when
$\varepsilon_0$ is bounded independently of $n$. The hypotheses
$0\le H_i\le1$ allow $\varepsilon_0$ of order $n$, and then
(\ref{eq:printed}) fails.

\section{A family with extensive ground energy}

Take qubits, $n\ge2$, and let $N_i=|1\rangle\langle1|$ on qubit $i$. Put
\begin{align*}
    H_i=\frac12+\frac14(N_i+N_{i+1}),\qquad 1\le i\le n-1.
\end{align*}
Each $H_i$ acts on the qubits $i$ and $i+1$ and has eigenvalues
$\tfrac12,\tfrac34,1$, so $0\le H_i\le1$. The terms commute, and
\begin{align*}
    H=\frac{n-1}2+\frac14(N_1+N_n)+\frac12\sum_{1<i<n}N_i
\end{align*}
is diagonal in the computational basis. Its unique ground state is
$|0\cdots0\rangle$, with $\varepsilon_0=(n-1)/2$, and its first excited states
flip one end qubit, so $\varepsilon=1/4$ for every $n$. Then
\begin{align*}
    C=\frac{2n}{n+1}\ge\frac43,
    \qquad
    \frac{1-\varepsilon_1/n}{1-\varepsilon_0/n}
    =\frac{n-\varepsilon_1}{n-\varepsilon_0}=1-\frac1{2n+2}.
\end{align*}
The choice of $m$ must make the bound of equation \texttt{e:cutdown} (line 603)
at most $c_\varepsilon/(2q(n))<1$. Since $-\log(1-x)\le x/(1-x)$, this forces
$m\ge(2n+1)\log(2q(n)/c_\varepsilon)$, and then
\begin{align*}
    C^{2m}\ge\Bigl(\frac43\Bigr)^{2(2n+1)\log(2q(n)/c_\varepsilon)}
    =\Bigl(\frac{2q(n)}{c_\varepsilon}\Bigr)^{2(2n+1)\log(4/3)},
\end{align*}
whose exponent grows linearly in $n$ at fixed $\varepsilon$. So
(\ref{eq:printed}) is false for this family, and the Bernstein estimate of the
proof, which needs $\ell$ of order $C^{2m}q(n)^2\log(d^n n^3)$, does not give a
polynomial number of words. Rescaling $K$, which the source allows (line 614),
does not help: the estimate controls $\|K-A\|$ relative to $\|A\|=1$, and the
ratio $R/\|A\|$ is unchanged by a common rescaling.

This note does not decide whether some other argument gives a polynomial $\ell$
for this family; it records only that the printed step and the stated
dependence of $\ell$ and of the running time on $\varepsilon$ alone do not
follow from the source's argument.

\section{The corrected statement}

With $R=C^m+1$, an accuracy $\eta>0$, and
\begin{align*}
    \ell\ge\frac{2R^2+2R\eta/3}{\eta^2}\log(8d^nn^3),\qquad
    \kappa\ge\frac{2m}n+2\log(4\ell n^4),
\end{align*}
the matrix Bernstein inequality, the scalar Bernstein inequality for the
number of occurrences of a letter, and a union bound over the $\ell$ words and
the $n$ letters give $\|K-A\|<\eta$ and at most $\kappa$ occurrences of every
letter in every word with probability at least $1-1/(2n^3)$. This holds for
every $m$ and every $\varepsilon_0$. At
$m=\lceil((n-\varepsilon_0)/\varepsilon)\log(1/\eta)\rceil$ and
$0<\eta\le1$, the bound $\varepsilon_0\le n-1$ gives
\begin{align*}
    C^m\le n\,\eta^{-\varepsilon_0/\varepsilon}.
\end{align*}
Hence $\ell$ is polynomial in $n$ and $1/\eta$ whenever $\varepsilon_0/\varepsilon$
is bounded, and $\kappa=O(\varepsilon^{-1}\log n)$ for $\eta=1/\mathrm{poly}(n)$
in that case, as in the source. For general $\varepsilon_0$, $\log\ell$ and
$\kappa$ carry the additional term
$O((\varepsilon_0/\varepsilon)\log(1/\eta))$, and the exponent of the running
time of Theorem~1 acquires a term of order
$\varepsilon_0/\varepsilon$.

\section{Formalization}

The sampling operator is \leanid{TNLean.MPS.AreaLaw.sampledProjector}, with
the explicit $\ell$ and $\kappa$ of the previous section
\leanid{TNLean.MPS.AreaLaw.sampledProjectorWordCount} and
\leanid{TNLean.MPS.AreaLaw.sampledProjectorLetterBound}. The probability form is
\leanid{TNLean.MPS.AreaLaw.one_sub_le_measureReal_sampledProjector}, and the
existence form, with the rank bound of equation \texttt{e:decomp}, is
\leanid{TNLean.MPS.AreaLaw.exists_sampledProjector}. The bounds on $C^m$ are
\leanid{TNLean.MPS.AreaLaw.powerNormalization_pow_le_exp} and
\leanid{TNLean.MPS.AreaLaw.powerNormalization_pow_ceil_le}. These statements
carry only the source's hypotheses. The statement $\ell=n^{O(1/\varepsilon)}$
for unbounded $\varepsilon_0$, and the running time of Theorem~1,
are not formalized; the status stays open until the dependence of the whole
algorithm on $\varepsilon_0$ is settled.

\bibliographystyle{unsrt}
\bibliography{references}
\end{document}
